\section{Descripción del problema}

\subsection{Contexto general}

La empresa distribuidora de combustibles opera de manera centralizada desde un deposito principal para el abastecimiento eficiente a cuatro estaciones de servicio. Para responder a la demanda de combustible Regular y Diésel y para ello la empresa dispone de dos camiones cisterna con tecnología de doble compartimiento. Esta configuración permite segregar fisicamente los combustibles, garantizando que cada espacio transporte únicamente un tipo de producto por compartimiento y eliminando cualquier riesgo de contaminación cruzada.

La gestión operativa enfrenta como desafío coordinar múltiples variables logísticas para minimizar el costo total de operación. Se deben diseñar rutas optimas las cuales minimicen las distancias y los costos fijos de la flota, respetando de manera estricta las ventanas de tiempo de las estaciones y evitando penalizaciones por demanda no satisfecha. Adicionalmente, el modelo integra restricciones complejas como el proceso de carga secuencial en el deposito y el control de la estabilidad de carga de los camiones en ruta.

\subsection{Datos del problema}

A continuación se presentan los datos principales del problema.

\begin{table}[H]
\centering
\caption{Demanda de combustible y ventanas de tiempo por estación}
\begin{tabular}{cccc}
\toprule
Estación & Regular (L) & Diésel (L) & Ventana de Tiempo \\
\midrule
1 & 3000 & 2000 & [06:00 - 10:00] \\
2 & 4000 & 1500 & [07:00 - 12:00] \\
3 & 2500 & 3000 & [08:00 - 14:00] \\
4 & 1000 & 2500 & [06:00 - 11:00] \\
\bottomrule
\end{tabular}
\end{table}
\begin{table}[H]
\centering
\caption{Características de la flota}
\begin{tabular}{cccccc}
\toprule
Camión & Cap. C0 (L) & Cap. C1 (L) & Costo fijo & Velocidad & Hora salida \\
\midrule
T1 & 8000 & 7000 & 500 & 60 km/h & 05:00 \\
T2 & 6000 & 9000 & 400 & 60 km/h & 05:30 \\
\bottomrule
\end{tabular}
\end{table}
